\chapter{Convertidor de impedancia negativa}

Es un circuito que se comporta como una resistencia de valor negativo \(Z=-R\). Se comporta como una carga que, en lugar de disipar energía, aporta energía al circuito. El diagrama topológico del circuito se muestra en la figura \ref{fig_nic}.
\begin{figure}[!ht]
  \centering
  \begin{circuitikz}
    \node[op amp,yscale=-1] (A) at (0,0) {};
    \draw (A.+) to[short] ++(0,1) coordinate(vin) to[R,l=$R$] (vin -| A.out) to[short,-*] (A.out);
    \draw (A.-) to[short] ++(0,-1) coordinate(grnd) to[R,l=$R$] (grnd -| A.out) to[short] (A.out);
    \draw (grnd) to[R,l_={$R$},*-] ++(-3,0) coordinate(len) node[ground]{};
    \draw (vin) to[short,*-o] (vin -| len) node[left]{$V_\text{in}$};
    \draw (A.out) to[short,-o] ++(0.5,0) node[right]{$V_\text{out}$};
  \end{circuitikz}
  \caption{Diagrama topológico del NIC.}
  \label{fig_nic}
\end{figure}
La impedancia equivalente es 
\[
  Z_\text{eq}=-R
\]

\section{Derivación}

Al NIC se le aplica un voltaje \(V_\text{in}\) y se genera una corriente \(I\). La impedancia equivalente puede ser calculada como 
\[
  Z_\text{eq} = \frac{V}{I}
\]
\begin{figure}[!ht]
  \centering
  \begin{circuitikz}
    \node[op amp,yscale=-1] (A) at (0,0) {};
    \draw (A.+) to[short] ++(0,1) coordinate(vin) to[R,l=$R$] (vin -| A.out) to[short,i={$I$},-*] (A.out);
    \draw (A.-) to[short] ++(0,-1) coordinate(grnd) to[R,l=$R$] (grnd -| A.out) to[short,i_<={$I_a$}] (A.out);
    \draw (grnd) to[R,l_={$R$},i={$I_b$},*-] ++(-3,0) coordinate(len) node[ground]{};
    \draw (vin) to[short,*-o] (vin -| len) node[left]{$V$};
    \draw (A.out) to[short,-o] ++(0.5,0) node[right]{$V_\text{out}$};
  \end{circuitikz}
  \caption{Análisis de corrientes del NIC.}
\end{figure}

En el nodo del terminal no inversor, se cumple
\begin{align}
  I_a&=I_b\notag\\ 
  \frac{V_\text{out}-V}{R}&=\frac{V-0}{R} \notag\\ 
  V_\text{out}-V &= V \notag \\ 
  V_\text{out} &= 2V \label{eq_vout_nic}
\end{align}
Calculamos la corriente en la rama superior
\begin{equation}\label{eq_nic_rama_sup}
  I=\frac{V-V_\text{out}}{R}
\end{equation}
Reemplazando \eqref{eq_vout_nic} en \eqref{eq_nic_rama_sup}:
\begin{align*}
  I&=\frac{V-2V}{R} \\ 
  I&=-\frac{V}{R} \\ 
  \frac{V}{I}&=-R \\ 
  Z_\text{eq}&= -R
\end{align*}

\section{Ejercicios para el lector}

\subsection{Preguntas teóricas}

\begin{enumerate}
  \item ¿Qué significa NIC?
  \item ¿Para qué sirve?
  \item ¿De donde obtiene la energía adicional el circuito NIC?
  \item ¿Cuántos amplificadores operacionales necesito para construir un NIC?
\end{enumerate}
